\chapter{PENDAHULUAN}
\section{Latar Belakang}


Tuliskan di sini \textbf{Pedoman Penulisan Latar Belakang:} Jelaskan alasan mendasar mengapa Anda memilih aplikasi mobile tersebut untuk dianalisis. Deskripsikan secara umum fungsi utama aplikasi, segmentasi pengguna, serta kontribusi aplikasi dalam memecahkan masalah komputasi mobile.


\section{Tujuan dan Manfaat}


Tuliskan di sini \textbf{Pedoman Penulisan Tujuan dan Manfaat:} Jelaskan tujuan dari pelaksanaan analisis aplikasi mobile ini (seperti mengidentifikasi arsitektur program dan menganalisis alur kerja). Sebutkan manfaat praktisnya bagi pembaca, institusi pendidikan, dan developer aplikasi terkait.


\section{Metode Penelitian}


Tuliskan di sini \textbf{Pedoman Penulisan Metode Penelitian:} Deskripsikan metode atau teknik pengumpulan data yang Anda gunakan (seperti melakukan pengujian mandiri terhadap aplikasi mobile, studi literatur pustaka mengenai platform Android/iOS, atau survei umpan balik pengguna).


\section{Ruang Lingkup}


Tuliskan di sini \textbf{Pedoman Penulisan Ruang Lingkup:} Batasi cakupan analisis program berbasis mobile ini agar pembahasan terfokus (misalnya membatasi pada analisis fitur transaksi utama, modul autentikasi user, atau performa alokasi memori pada versi sistem operasi tertentu).

